
\documentclass{article}
\usepackage{colortbl}
\usepackage{makecell}
\usepackage{multirow}
\usepackage{supertabular}

\begin{document}

\newcounter{utterance}

\centering \large Interaction Transcript for game `bbh', experiment `multistep\_arithmetic\_two', episode 94 with reflexion{-}Qwen3.5{-}27B.
\vspace{24pt}

{ \footnotesize  \setcounter{utterance}{1}
\setlength{\tabcolsep}{0pt}
\begin{supertabular}{c@{$\;$}|p{.15\linewidth}@{}p{.15\linewidth}p{.15\linewidth}p{.15\linewidth}p{.15\linewidth}p{.15\linewidth}}
    \# & \multicolumn{2}{c}{Player} && \multicolumn{2}{c}{Game Master} \\
    \hline

    \theutterance \stepcounter{utterance}  
    & & & \multicolumn{4}{p{0.6\linewidth}}{
        \cellcolor[rgb]{0.9,0.9,0.9}{
            \makecell[{{p{\linewidth}}}]{
                \texttt{\tiny{[P1$\langle$GM]}}
                \texttt{You have attempted this same task before and did not succeed. These are the plans you wrote after those attempts.} \\
\texttt{Trial \#1: New plan: My previous attempt failed because I provided a verbose, step{-}by{-}step explanation instead of the single numerical answer required by the pattern established in the previous examples (Q: ... A: {[}number{]}). The game expects only the final result. Next time, I will calculate the expression internally using the order of operations but output **only** the final integer, omitting all intermediate steps, reasoning, and formatting like "Step 1" or "Final Answer". For the current problem, the calculation is \$({-}7 + 180) \textbackslash{}times ({-}35 {-} 6 {-} 7) = 173 \textbackslash{}times {-}48 = {-}8304\$, so I will simply output "{-}8304".} \\
\\ 
\\ 
\texttt{Solve multi{-}step arithmetic problems.} \\
\\ 
\texttt{Q: (({-}5 + 9 * {-}4 {-} 0) * (4 + {-}7 + 0 * {-}5)) =} \\
\texttt{A: 123} \\
\\ 
\texttt{Q: (({-}9 * 7 * 7 * {-}9) + (4 * {-}9 {-} 8 {-} {-}4)) =} \\
\texttt{A: 3929} \\
\\ 
\texttt{Q: (({-}3 + 5 * 8 * {-}4) {-} (9 {-} 8 * {-}7 + {-}9)) =} \\
\texttt{A: {-}219} \\
\\ 
\texttt{Q: (({-}7 + 5 * {-}6 * {-}6) * (7 * {-}5 + {-}6 {-} 7)) =} \\
\texttt{A:} \\
            }
        }
    }
     \\ \\

    \theutterance \stepcounter{utterance}  
    & \multicolumn{4}{p{0.6\linewidth}}{
        \cellcolor[rgb]{0.561,0.737,0.561}{
            \makecell[{{p{\linewidth}}}]{
                \texttt{\tiny{[P1$\rangle$GM]}}
                \texttt{{-}1080} \\
            }
        }
    }
    & & \\ \\

    \theutterance \stepcounter{utterance}  
    & & & \multicolumn{2}{p{0.3\linewidth}}{
        \cellcolor[rgb]{0.9,0.9,0.9}{
            \makecell[{{p{\linewidth}}}]{
                \texttt{\tiny{[GM$|$GM]}}
                \texttt{Target: {-}8304} \\
            }
        }
    }
    & & \\ \\

    \theutterance \stepcounter{utterance}  
    & & & \multicolumn{2}{p{0.3\linewidth}}{
        \cellcolor[rgb]{0.9,0.9,0.9}{
            \makecell[{{p{\linewidth}}}]{
                \texttt{\tiny{[GM$|$GM]}}
                \texttt{game\_result = LOSE} \\
            }
        }
    }
    & & \\ \\

    \theutterance \stepcounter{utterance}  
    & & \multicolumn{4}{p{0.6\linewidth}}{
        \cellcolor[rgb]{0.95,0.95,0.95}{
            \makecell[{{p{\linewidth}}}]{
                \texttt{\tiny{[P1$\leftrightarrow$P]}}
                \texttt{Post{-}episode Reflexion (not sent to GM):} \\
\texttt{New plan: My previous attempt failed because I miscalculated the arithmetic expression, likely due to an error in handling the order of operations or sign changes, resulting in an incorrect answer ({-}1080 instead of the correct {-}8304). The strategy of providing only the final number was correct, but the execution of the math was flawed. Next time, I will rigorously break down the expression into smaller parts before combining them: first, calculate the inner parentheses strictly following PEMDAS (multiplication before addition/subtraction), paying close attention to negative signs. Specifically, for the first part \$({-}7 + 5 \textbackslash{}times {-}6 \textbackslash{}times {-}6)\$, I will compute \$5 \textbackslash{}times {-}6 = {-}30\$, then \${-}30 \textbackslash{}times {-}6 = 180\$, and finally \${-}7 + 180 = 173\$. For the second part \$(7 \textbackslash{}times {-}5 + {-}6 {-} 7)\$, I will compute \$7 \textbackslash{}times {-}5 = {-}35\$, then \${-}35 {-} 6 {-} 7 = {-}48\$. Finally, I will multiply the two results: \$173 \textbackslash{}times {-}48 = {-}8304\$. I will output only this verified final integer.} \\
            }
        }
    }
    & \\ \\

\end{supertabular}
}

\end{document}
